\begin{theorem}[Fixed-calibration stabilization of reusable menus]
\label{thm:v30-menu-finite}
Fix a finite graph and a calibration $a$. For every $K\ge1$ and every
nonempty $\mathcal E\subset(0,1]$,
\begin{equation}\label{eq:v30-menu-finite}
 0\le r_K(\mathcal E,a)\le r_1((0,1],a)<\infty.
\end{equation}
There are an order $\pi_*$ and a finite $X\ge0$ such that
\begin{equation}\label{eq:v30-eventual-order}
 W_{\pi_*}(2^{-x},a)=W_G(2^{-x},a)\qquad(x\ge X).
\end{equation}
For a prescribed nonempty menu $\Pi$, put
\[
 d_e=\max\{0\le\ell\le p_e:V_{e,\ell}(a)>0\},\qquad
 D_\pi=\frac12\max_j\sum_{e\in\delta(S_j^\pi)}d_e,\qquad
 D_* =\min_\pi D_\pi.
\]
The prescribed menu has finite excess on $(0,1]$ if and only if
$\min_{\pi\in\Pi}D_\pi=D_*$. These statements include exact
collisions and the empty graph.
\end{theorem}
\begin{proof}
Omit all zero-volume branches and write
\[
 b_e(2^{-x},a)=\max_{\ell:V_{e,\ell}>0}
                 \{(\ell/2)x+c_{e,\ell}\},\qquad
 c_{e,\ell}=\log_2V_{e,\ell}(a),\quad c_{e,0}=0.
\]
Every retained intercept is finite. If $d_e>0$, all sufficiently
large $x$ satisfy
\[
 b_e(2^{-x},a)=(d_e/2)x+c_{e,d_e};
\]
indeed it suffices that
$x\ge\max_{\ell<d_e}[2(c_{e,\ell}-c_{e,d_e})/(d_e-\ell)]_+$.
An edge with only its zero branch contributes zero.

A cut sum is a maximum of finitely many affine functions, by expanding
the independent finite maxima. The same is true of $W_\pi$.
Beyond a finite common threshold, every $W_\pi$ is consequently of
the form $D_\pi x+C_\pi$. Its slope is the displayed maximum cut
sum of the $d_e/2$; among cuts with maximal slope its intercept is
the largest sum of the edge intercepts. Since there are only finitely
many orders, their minimum is eventually the affine function with
lexicographically least pair $(D_\pi,C_\pi)$. Any order realizing
that pair gives \eqref{eq:v30-eventual-order} after increasing the
threshold to include their finitely many crossings.

On $[0,X]$, all order profiles and their finite minimum are continuous.
Thus $W_{\pi_*}-W_G$ is bounded there and is zero thereafter. The
singleton menu $\{\pi_*\}$ proves \eqref{eq:v30-menu-finite}.
Every menu's excess is nonnegative. A prescribed menu is eventually
affine with slope $\min_{\pi\in\Pi}D_\pi$; its eventual excess has
that slope minus $D_*$. The difference is nonnegative and is bounded
on the half-line precisely when its slope is zero. For the empty graph
all profiles, slopes and excesses are zero.
\end{proof}

\begin{proposition}[A finite critical-resolution formula]
\label{prop:v30-critical-resolutions}
At a fixed calibration there is a finite set $T(a)\subset[0,\infty)$,
containing zero and depending only on the graph's affine branches,
such that
\begin{equation}\label{eq:v30-critical-resolutions}
 r_K((0,1],a)=
 \min_{\substack{1\le|\Pi|\le K\\\min_{\pi\in\Pi}D_\pi=D_*}}
 \max_{x\in T(a)}
       \{W_\Pi(2^{-x},a)-W_G(2^{-x},a)\}.
\end{equation}
Thus the full fixed-calibration, all-resolution excess is a finite
minimax problem, even though its accuracy domain is not bounded away
from zero.
\end{proposition}
\begin{proof}
Let $\mathcal A$ consist of the affine functions obtained by choosing
one retained edge branch on every edge of each cut of each order and
summing them. Include the zero function. The set is finite, and
$W_\pi$ is a maximum over a subset of $\mathcal A$.
Take $T(a)$ to be zero together with all nonnegative intersections of
two nonparallel functions in $\mathcal A$. Repeated equal functions
cause no additional intersections. Between two consecutive points of
$T(a)$, and on its final unbounded interval, the ordering of all affine
functions is fixed. Every finite maximum and minimum made from them
therefore agrees with one affine function on that interval.

For a menu with $\min_{\pi\in\Pi}D_\pi=D_*$, the excess is constant
on the final interval and affine on every preceding interval. Its
supremum is hence attained at a point of $T(a)$. Menus failing that
slope condition have infinite excess by
Theorem~\ref{thm:v30-menu-finite} and cannot minimize, since a finite
competitor exists. Taking the minimum over the remaining finite family
proves \eqref{eq:v30-critical-resolutions}. The construction is an
exact finite characterization with the given real intercepts, not a
polynomial-time algorithm or a numerical enclosure of unknown
coefficients.
\end{proof}

\begin{corollary}[Uniformity away from vanishing retained volumes]
\label{cor:v30-stratum-uniform}
Suppose a calibration set has the same retained index sets on every
edge, and all their positive volumes are bounded above and below by
positive constants. Then one deterministic order has bounded excess
uniformly on that entire calibration set and on all of $(0,1]$.
\end{corollary}
\begin{proof}
The integers $d_e$ are constant and the retained intercepts are
uniformly bounded. The edge stabilization thresholds in the preceding
proof can therefore be bounded uniformly. The eventual cut slopes are
also fixed, and their distinct values differ by at least $1/2$.
Their intercepts are uniformly bounded finite sums. Consequently
there is a common threshold beyond which every order has its fixed
slope $D_\pi$ and a uniformly bounded intercept. Choose any order
minimizing $D_\pi$; its excess is bounded uniformly on that final
interval. On the common bounded initial interval, uniform boundedness
of the edge branches gives the same conclusion. This does not require
the minimizing eventual intercept to be realized by the same order
at all calibrations.
\end{proof}

The optimized functional at one fixed calibration is therefore always
finite. A prescribed menu of larger eventual slope can have infinite
excess. A bound uniform over calibrations approaching a collision can
also diverge, because retained volumes and their intercepts need not
have the lower bounds in Corollary~\ref{cor:v30-stratum-uniform}.
The collision-path results below concern this latter quantifier and
are unchanged.
